%==============================================================
%  estilos/codigo.tex
%  Utilidades para mostrar código y pseudocódigo en la tesis.
%
%  Motor: listings (resaltado de sintaxis) + tcolorbox (marco y título).
%  No requiere instalar nada ni compilar con --shell-escape: funciona con
%  pdflatex/latexmk tal como está configurada la tesis.
%
%  Se carga desde el preámbulo de tesis.tex con  %==============================================================
%  estilos/codigo.tex
%  Utilidades para mostrar código y pseudocódigo en la tesis.
%
%  Motor: listings (resaltado de sintaxis) + tcolorbox (marco y título).
%  No requiere instalar nada ni compilar con --shell-escape: funciona con
%  pdflatex/latexmk tal como está configurada la tesis.
%
%  Se carga desde el preámbulo de tesis.tex con  %==============================================================
%  estilos/codigo.tex
%  Utilidades para mostrar código y pseudocódigo en la tesis.
%
%  Motor: listings (resaltado de sintaxis) + tcolorbox (marco y título).
%  No requiere instalar nada ni compilar con --shell-escape: funciona con
%  pdflatex/latexmk tal como está configurada la tesis.
%
%  Se carga desde el preámbulo de tesis.tex con  %==============================================================
%  estilos/codigo.tex
%  Utilidades para mostrar código y pseudocódigo en la tesis.
%
%  Motor: listings (resaltado de sintaxis) + tcolorbox (marco y título).
%  No requiere instalar nada ni compilar con --shell-escape: funciona con
%  pdflatex/latexmk tal como está configurada la tesis.
%
%  Se carga desde el preámbulo de tesis.tex con  \input{estilos/codigo}
%  Documentación y ejemplos: notes/codigo.md  y  notes/codigo-ejemplos.tex
%
%  Resumen de la interfaz:
%    \begin{codigo}[style=python]{Título}   ... \end{codigo}          (se parte entre páginas)
%    \begin{codigoflotante}[style=python]{Título} ... \end{codigoflotante}  (flotante)
%    \archivocodigo[firstline=1,lastline=20]{ruta.py}{Título}         (desde archivo)
%    \py{código}   \sh{comando}   \cod{texto}                         (en línea)
%==============================================================

\usepackage{xcolor}          % Modelos de color RGB (ya lo carga tcolorbox, aquí para dejarlo explícito)
\usepackage{listings}        % Resaltado de sintaxis
\usepackage{upquote}         % Comillas rectas ' en el código (no tipográficas)
\usepackage{listingsutf8}    % Archivos completos en UTF-8 con \lstinputlisting
\usepackage[listings,breakable,skins,hooks]{tcolorbox}   % Marco, título, partido entre páginas

%--------------------------------------------------------------
%  PALETA
%  Sobria a propósito: la tesis se imprime en blanco y negro, así que
%  todo debe seguir siendo distinguible por contraste y por estilo.
%--------------------------------------------------------------
\definecolor{codigoFondo}{RGB}{252,252,250}      % fondo del listado
\definecolor{codigoTitulo}{RGB}{238,238,232}     % barra del título
\definecolor{codigoLinea}{RGB}{130,130,130}      % marco
\definecolor{codigoGris}{RGB}{95,95,95}          % números de línea y flecha de continuación
\definecolor{codigoClave}{RGB}{0,55,125}         % palabras clave
\definecolor{codigoCadena}{RGB}{140,35,35}       % cadenas
\definecolor{codigoComentario}{RGB}{85,105,85}   % comentarios
\definecolor{codigoTipo}{RGB}{20,85,105}         % tipos y clases

%--------------------------------------------------------------
%  CONFIGURACIÓN COMÚN
%  El mapa «literate» es lo que permite escribir acentos, eñes y
%  comillas latinas DENTRO de un listado con pdflatex: listings no
%  entiende UTF-8 multibyte por sí solo.
%--------------------------------------------------------------
\lstset{
  basicstyle=\ttfamily\small,
  columns=fullflexible,
  keepspaces=true,
  showstringspaces=false,
  tabsize=4,
  upquote=true,
  breaklines=true,                              % parte las líneas largas
  breakatwhitespace=false,
  postbreak=\mbox{\textcolor{codigoGris}{$\hookrightarrow$}\space},
  numbers=left,
  numberstyle=\tiny\color{codigoGris},
  numbersep=8pt,
  xleftmargin=17pt,
  keywordstyle=\color{codigoClave}\bfseries,
  stringstyle=\color{codigoCadena},
  commentstyle=\color{codigoComentario}\itshape,
  emphstyle=\color{codigoTipo},
  captionpos=b,
  literate=%
    {á}{{\'a}}1 {é}{{\'e}}1 {í}{{\'i}}1 {ó}{{\'o}}1 {ú}{{\'u}}1
    {Á}{{\'A}}1 {É}{{\'E}}1 {Í}{{\'I}}1 {Ó}{{\'O}}1 {Ú}{{\'U}}1
    {ü}{{\"u}}1 {Ü}{{\"U}}1 {ñ}{{\~n}}1 {Ñ}{{\~N}}1
    {¿}{{?`}}1 {¡}{{!`}}1
    {«}{{\guillemotleft}}1 {»}{{\guillemotright}}1
    {—}{{---}}1 {–}{{--}}1 {…}{{\ldots}}1
    {º}{{\textordmasculine}}1 {ª}{{\textordfeminine}}1 {°}{{\textdegree}}1
    {→}{{$\rightarrow$}}1 {≈}{{$\approx$}}1 {∈}{{$\in$}}1
}

%--------------------------------------------------------------
%  LENGUAJES
%  Python, bash y SQL vienen con listings; JSON y el pseudocódigo
%  en español se definen aquí porque listings no los trae.
%--------------------------------------------------------------
\lstdefinestyle{python}{language=Python}

\lstdefinestyle{bash}{language=bash}

\lstdefinestyle{sql}{language=SQL}

\lstdefinelanguage{json}{
  alsoletter={-},
  morestring=[b]",
  morekeywords={true,false,null},
  sensitive=true,
  morecomment=[l]{//},
}

\lstdefinestyle{json}{language=json}

\lstdefinelanguage{seudo}{
  sensitive=false,
  morekeywords={algoritmo,entrada,salida,para,cada,hacer,fin,mientras,si,entonces,
    sino,devolver,retornar,repetir,hasta,funcion,procedimiento,inicio,
    verdadero,falso,no,y,o,listas,consulta,consulta,documento,documentos,
    ordenar,mezclar,fusionar,rango,puntuacion,puntuaciones},
  morecomment=[l]{\#},
}

\lstdefinestyle{seudocodigo}{language=seudo}

% Salida de consola: sin números de línea ni resaltado, para pegarla tal cual.
\lstdefinestyle{salida}{
  numbers=none,
  xleftmargin=4pt,
  keywordstyle={},
  stringstyle={},
  commentstyle=\color{codigoGris},
}

%--------------------------------------------------------------
%  CAJAS (tcolorbox)
%--------------------------------------------------------------
\tcbset{
  codigoCaja/.style={
    listing engine=listings,
    listing only,
    colback=codigoFondo,
    colframe=codigoLinea,
    coltitle=black,
    colbacktitle=codigoTitulo,
    fonttitle=\bfseries\small,
    boxrule=0.4pt,
    arc=1.5pt,
    outer arc=1.5pt,
    left=4pt, right=4pt, top=4pt, bottom=4pt,
    before skip=8pt, after skip=8pt,
  },
  % Alias cortos: \begin{codigo}[py]{…} equivale a [listing style=python].
  % La forma canónica de tcolorbox, [listing style=python] o
  % [listing options={…}], también funciona y sirve para cualquier opción
  % de listings (numbers=none, firstnumber=5, …).
  py/.style={listing style=python},
  sh/.style={listing style=bash},
  sql/.style={listing style=sql},
  js/.style={listing style=json},
  seudo/.style={listing style=seudocodigo},
  sal/.style={listing style=salida},
  % Rango de líneas de un archivo: \archivocodigo[lineas={10}{25}]{ruta}{Título}
  lineas/.style 2 args={listing options app={firstline=#1, lastline=#2}},
}

%--------------------------------------------------------------
%  ENTORNOS NUMERADOS
%  Los tres comparten el contador «listadocodigo», que se reinicia en
%  cada capítulo: el título sale como «Listado 4.1 · …».
%--------------------------------------------------------------
\newtcblisting[auto counter, number within=chapter, list inside=lol, list type=lstlisting]%
  {codigo}[2][]{%
  codigoCaja,
  breakable,                                   % se puede partir entre páginas
  list text={#2},                              % el «Índice de listados» muestra número y título
  title={Listado~\thetcbcounter\ \textperiodcentered\ #2},
  #1
}

\newtcblisting[use counter from=codigo, list inside=lol, list type=lstlisting]%
  {codigoflotante}[2][]{%
  codigoCaja,
  float, floatplacement=tbp,
  list text={#2},
  title={Listado~\thetcbcounter\ \textperiodcentered\ #2},
  #1
}

\newtcbinputlisting[use counter from=codigo, list inside=lol, list type=lstlisting]%
  {\archivocodigo}[3][]{%
  codigoCaja,
  breakable,
  list text={#3},
  title={Listado~\thetcbcounter\ \textperiodcentered\ #3},
  listing file={#2},
  #1
}

%--------------------------------------------------------------
%  CÓDIGO EN LÍNEA
%  Se definen SIN argumento a propósito: \lstinline lee el texto en modo
%  verbatim del flujo de entrada, y con el código como argumento de macro
%  falla («lstinline ended by EOL»). El delimitador es el carácter que
%  sigue a la llamada:
%      \py{revision="1.2"}     \sh|ssh iimas|     \cod{top-100}
%  Ojo: el contenido no puede llevar el carácter delimitador; para casos
%  raros, usa \verb|...| .
%--------------------------------------------------------------
\newcommand{\py}{\lstinline[style=python]}
\newcommand{\sh}{\lstinline[style=bash]}
\newcommand{\cod}{\lstinline}

%--------------------------------------------------------------
%  NOMBRES EN ESPAÑOL
%  Se fijan en \AtBeginDocument para que manden sobre los de babel.
%--------------------------------------------------------------
\AtBeginDocument{%
  \renewcommand{\lstlistingname}{Listado}%
  \renewcommand{\lstlistlistingname}{Índice de listados}%
}

%--------------------------------------------------------------
%  REFERENCIAS
%  Cualquiera de los tres entornos acepta la clave «label»:
%      \begin{codigoflotante}[py, label={lst:carga}]{Carga del split}
%  Así, \ref{lst:carga} imprime «5.1»; \autoref{lst:carga} (hyperref)
%  imprime «Listado 5.1»; y \cref{lst:carga} (cleveref) «listado 5.1».
%  cleveref es opcional: si no se carga, nada de esto estorba.
%  Dos cosas que costaron encontrar y conviene no volver a romper:
%   - El contador que crea tcolorbox se llama tcb@cnt@codigo. Es el nombre
%     con el que hay que registrar los nombres en español.
%   - Hay que registrarlos AL CARGAR cleveref, no en \AtBeginDocument:
%     cleveref construye los formatos de referencia al terminar de cargarse,
%     de modo que una definición posterior llega tarde y \cref imprime «??»
%     (aviso «reference format for label type … undefined»). Para eso está
%     el gancho \AddToHook{package/cleveref/after} de abajo.
%--------------------------------------------------------------
\AddToHook{package/cleveref/after}{%
  \Crefname{tcb@cnt@codigo}{Listado}{Listados}%
  \crefname{tcb@cnt@codigo}{listado}{listados}%
}

\makeatletter
\AtBeginDocument{%
  \expandafter\newcommand\csname tcb@cnt@codigoautorefname\endcsname{Listado}%
  %--- ÍNDICE DE LISTADOS EN EL ÍNDICE GENERAL
  % tocbibind solo conoce las listas de figuras y de tablas, así que el
  % «Índice de listados» no aparecía en el índice general. Se reusa su
  % propio mecanismo (\tocfile), que pone el encabezado y DESPUÉS el
  % \addcontentsline: así la entrada apunta a la página del encabezado
  % aunque la lista ocupe varias páginas.
  \@ifpackageloaded{tocbibind}{%
    \renewcommand{\lstlistoflistings}{\tocfile{\lstlistlistingname}{lol}}%
  }{}%
}
\makeatother

%  Documentación y ejemplos: notes/codigo.md  y  notes/codigo-ejemplos.tex
%
%  Resumen de la interfaz:
%    \begin{codigo}[style=python]{Título}   ... \end{codigo}          (se parte entre páginas)
%    \begin{codigoflotante}[style=python]{Título} ... \end{codigoflotante}  (flotante)
%    \archivocodigo[firstline=1,lastline=20]{ruta.py}{Título}         (desde archivo)
%    \py{código}   \sh{comando}   \cod{texto}                         (en línea)
%==============================================================

\usepackage{xcolor}          % Modelos de color RGB (ya lo carga tcolorbox, aquí para dejarlo explícito)
\usepackage{listings}        % Resaltado de sintaxis
\usepackage{upquote}         % Comillas rectas ' en el código (no tipográficas)
\usepackage{listingsutf8}    % Archivos completos en UTF-8 con \lstinputlisting
\usepackage[listings,breakable,skins,hooks]{tcolorbox}   % Marco, título, partido entre páginas

%--------------------------------------------------------------
%  PALETA
%  Sobria a propósito: la tesis se imprime en blanco y negro, así que
%  todo debe seguir siendo distinguible por contraste y por estilo.
%--------------------------------------------------------------
\definecolor{codigoFondo}{RGB}{252,252,250}      % fondo del listado
\definecolor{codigoTitulo}{RGB}{238,238,232}     % barra del título
\definecolor{codigoLinea}{RGB}{130,130,130}      % marco
\definecolor{codigoGris}{RGB}{95,95,95}          % números de línea y flecha de continuación
\definecolor{codigoClave}{RGB}{0,55,125}         % palabras clave
\definecolor{codigoCadena}{RGB}{140,35,35}       % cadenas
\definecolor{codigoComentario}{RGB}{85,105,85}   % comentarios
\definecolor{codigoTipo}{RGB}{20,85,105}         % tipos y clases

%--------------------------------------------------------------
%  CONFIGURACIÓN COMÚN
%  El mapa «literate» es lo que permite escribir acentos, eñes y
%  comillas latinas DENTRO de un listado con pdflatex: listings no
%  entiende UTF-8 multibyte por sí solo.
%--------------------------------------------------------------
\lstset{
  basicstyle=\ttfamily\small,
  columns=fullflexible,
  keepspaces=true,
  showstringspaces=false,
  tabsize=4,
  upquote=true,
  breaklines=true,                              % parte las líneas largas
  breakatwhitespace=false,
  postbreak=\mbox{\textcolor{codigoGris}{$\hookrightarrow$}\space},
  numbers=left,
  numberstyle=\tiny\color{codigoGris},
  numbersep=8pt,
  xleftmargin=17pt,
  keywordstyle=\color{codigoClave}\bfseries,
  stringstyle=\color{codigoCadena},
  commentstyle=\color{codigoComentario}\itshape,
  emphstyle=\color{codigoTipo},
  captionpos=b,
  literate=%
    {á}{{\'a}}1 {é}{{\'e}}1 {í}{{\'i}}1 {ó}{{\'o}}1 {ú}{{\'u}}1
    {Á}{{\'A}}1 {É}{{\'E}}1 {Í}{{\'I}}1 {Ó}{{\'O}}1 {Ú}{{\'U}}1
    {ü}{{\"u}}1 {Ü}{{\"U}}1 {ñ}{{\~n}}1 {Ñ}{{\~N}}1
    {¿}{{?`}}1 {¡}{{!`}}1
    {«}{{\guillemotleft}}1 {»}{{\guillemotright}}1
    {—}{{---}}1 {–}{{--}}1 {…}{{\ldots}}1
    {º}{{\textordmasculine}}1 {ª}{{\textordfeminine}}1 {°}{{\textdegree}}1
    {→}{{$\rightarrow$}}1 {≈}{{$\approx$}}1 {∈}{{$\in$}}1
}

%--------------------------------------------------------------
%  LENGUAJES
%  Python, bash y SQL vienen con listings; JSON y el pseudocódigo
%  en español se definen aquí porque listings no los trae.
%--------------------------------------------------------------
\lstdefinestyle{python}{language=Python}

\lstdefinestyle{bash}{language=bash}

\lstdefinestyle{sql}{language=SQL}

\lstdefinelanguage{json}{
  alsoletter={-},
  morestring=[b]",
  morekeywords={true,false,null},
  sensitive=true,
  morecomment=[l]{//},
}

\lstdefinestyle{json}{language=json}

\lstdefinelanguage{seudo}{
  sensitive=false,
  morekeywords={algoritmo,entrada,salida,para,cada,hacer,fin,mientras,si,entonces,
    sino,devolver,retornar,repetir,hasta,funcion,procedimiento,inicio,
    verdadero,falso,no,y,o,listas,consulta,consulta,documento,documentos,
    ordenar,mezclar,fusionar,rango,puntuacion,puntuaciones},
  morecomment=[l]{\#},
}

\lstdefinestyle{seudocodigo}{language=seudo}

% Salida de consola: sin números de línea ni resaltado, para pegarla tal cual.
\lstdefinestyle{salida}{
  numbers=none,
  xleftmargin=4pt,
  keywordstyle={},
  stringstyle={},
  commentstyle=\color{codigoGris},
}

%--------------------------------------------------------------
%  CAJAS (tcolorbox)
%--------------------------------------------------------------
\tcbset{
  codigoCaja/.style={
    listing engine=listings,
    listing only,
    colback=codigoFondo,
    colframe=codigoLinea,
    coltitle=black,
    colbacktitle=codigoTitulo,
    fonttitle=\bfseries\small,
    boxrule=0.4pt,
    arc=1.5pt,
    outer arc=1.5pt,
    left=4pt, right=4pt, top=4pt, bottom=4pt,
    before skip=8pt, after skip=8pt,
  },
  % Alias cortos: \begin{codigo}[py]{…} equivale a [listing style=python].
  % La forma canónica de tcolorbox, [listing style=python] o
  % [listing options={…}], también funciona y sirve para cualquier opción
  % de listings (numbers=none, firstnumber=5, …).
  py/.style={listing style=python},
  sh/.style={listing style=bash},
  sql/.style={listing style=sql},
  js/.style={listing style=json},
  seudo/.style={listing style=seudocodigo},
  sal/.style={listing style=salida},
  % Rango de líneas de un archivo: \archivocodigo[lineas={10}{25}]{ruta}{Título}
  lineas/.style 2 args={listing options app={firstline=#1, lastline=#2}},
}

%--------------------------------------------------------------
%  ENTORNOS NUMERADOS
%  Los tres comparten el contador «listadocodigo», que se reinicia en
%  cada capítulo: el título sale como «Listado 4.1 · …».
%--------------------------------------------------------------
\newtcblisting[auto counter, number within=chapter, list inside=lol, list type=lstlisting]%
  {codigo}[2][]{%
  codigoCaja,
  breakable,                                   % se puede partir entre páginas
  list text={#2},                              % el «Índice de listados» muestra número y título
  title={Listado~\thetcbcounter\ \textperiodcentered\ #2},
  #1
}

\newtcblisting[use counter from=codigo, list inside=lol, list type=lstlisting]%
  {codigoflotante}[2][]{%
  codigoCaja,
  float, floatplacement=tbp,
  list text={#2},
  title={Listado~\thetcbcounter\ \textperiodcentered\ #2},
  #1
}

\newtcbinputlisting[use counter from=codigo, list inside=lol, list type=lstlisting]%
  {\archivocodigo}[3][]{%
  codigoCaja,
  breakable,
  list text={#3},
  title={Listado~\thetcbcounter\ \textperiodcentered\ #3},
  listing file={#2},
  #1
}

%--------------------------------------------------------------
%  CÓDIGO EN LÍNEA
%  Se definen SIN argumento a propósito: \lstinline lee el texto en modo
%  verbatim del flujo de entrada, y con el código como argumento de macro
%  falla («lstinline ended by EOL»). El delimitador es el carácter que
%  sigue a la llamada:
%      \py{revision="1.2"}     \sh|ssh iimas|     \cod{top-100}
%  Ojo: el contenido no puede llevar el carácter delimitador; para casos
%  raros, usa \verb|...| .
%--------------------------------------------------------------
\newcommand{\py}{\lstinline[style=python]}
\newcommand{\sh}{\lstinline[style=bash]}
\newcommand{\cod}{\lstinline}

%--------------------------------------------------------------
%  NOMBRES EN ESPAÑOL
%  Se fijan en \AtBeginDocument para que manden sobre los de babel.
%--------------------------------------------------------------
\AtBeginDocument{%
  \renewcommand{\lstlistingname}{Listado}%
  \renewcommand{\lstlistlistingname}{Índice de listados}%
}

%--------------------------------------------------------------
%  REFERENCIAS
%  Cualquiera de los tres entornos acepta la clave «label»:
%      \begin{codigoflotante}[py, label={lst:carga}]{Carga del split}
%  Así, \ref{lst:carga} imprime «5.1»; \autoref{lst:carga} (hyperref)
%  imprime «Listado 5.1»; y \cref{lst:carga} (cleveref) «listado 5.1».
%  cleveref es opcional: si no se carga, nada de esto estorba.
%  Dos cosas que costaron encontrar y conviene no volver a romper:
%   - El contador que crea tcolorbox se llama tcb@cnt@codigo. Es el nombre
%     con el que hay que registrar los nombres en español.
%   - Hay que registrarlos AL CARGAR cleveref, no en \AtBeginDocument:
%     cleveref construye los formatos de referencia al terminar de cargarse,
%     de modo que una definición posterior llega tarde y \cref imprime «??»
%     (aviso «reference format for label type … undefined»). Para eso está
%     el gancho \AddToHook{package/cleveref/after} de abajo.
%--------------------------------------------------------------
\AddToHook{package/cleveref/after}{%
  \Crefname{tcb@cnt@codigo}{Listado}{Listados}%
  \crefname{tcb@cnt@codigo}{listado}{listados}%
}

\makeatletter
\AtBeginDocument{%
  \expandafter\newcommand\csname tcb@cnt@codigoautorefname\endcsname{Listado}%
  %--- ÍNDICE DE LISTADOS EN EL ÍNDICE GENERAL
  % tocbibind solo conoce las listas de figuras y de tablas, así que el
  % «Índice de listados» no aparecía en el índice general. Se reusa su
  % propio mecanismo (\tocfile), que pone el encabezado y DESPUÉS el
  % \addcontentsline: así la entrada apunta a la página del encabezado
  % aunque la lista ocupe varias páginas.
  \@ifpackageloaded{tocbibind}{%
    \renewcommand{\lstlistoflistings}{\tocfile{\lstlistlistingname}{lol}}%
  }{}%
}
\makeatother

%  Documentación y ejemplos: notes/codigo.md  y  notes/codigo-ejemplos.tex
%
%  Resumen de la interfaz:
%    \begin{codigo}[style=python]{Título}   ... \end{codigo}          (se parte entre páginas)
%    \begin{codigoflotante}[style=python]{Título} ... \end{codigoflotante}  (flotante)
%    \archivocodigo[firstline=1,lastline=20]{ruta.py}{Título}         (desde archivo)
%    \py{código}   \sh{comando}   \cod{texto}                         (en línea)
%==============================================================

\usepackage{xcolor}          % Modelos de color RGB (ya lo carga tcolorbox, aquí para dejarlo explícito)
\usepackage{listings}        % Resaltado de sintaxis
\usepackage{upquote}         % Comillas rectas ' en el código (no tipográficas)
\usepackage{listingsutf8}    % Archivos completos en UTF-8 con \lstinputlisting
\usepackage[listings,breakable,skins,hooks]{tcolorbox}   % Marco, título, partido entre páginas

%--------------------------------------------------------------
%  PALETA
%  Sobria a propósito: la tesis se imprime en blanco y negro, así que
%  todo debe seguir siendo distinguible por contraste y por estilo.
%--------------------------------------------------------------
\definecolor{codigoFondo}{RGB}{252,252,250}      % fondo del listado
\definecolor{codigoTitulo}{RGB}{238,238,232}     % barra del título
\definecolor{codigoLinea}{RGB}{130,130,130}      % marco
\definecolor{codigoGris}{RGB}{95,95,95}          % números de línea y flecha de continuación
\definecolor{codigoClave}{RGB}{0,55,125}         % palabras clave
\definecolor{codigoCadena}{RGB}{140,35,35}       % cadenas
\definecolor{codigoComentario}{RGB}{85,105,85}   % comentarios
\definecolor{codigoTipo}{RGB}{20,85,105}         % tipos y clases

%--------------------------------------------------------------
%  CONFIGURACIÓN COMÚN
%  El mapa «literate» es lo que permite escribir acentos, eñes y
%  comillas latinas DENTRO de un listado con pdflatex: listings no
%  entiende UTF-8 multibyte por sí solo.
%--------------------------------------------------------------
\lstset{
  basicstyle=\ttfamily\small,
  columns=fullflexible,
  keepspaces=true,
  showstringspaces=false,
  tabsize=4,
  upquote=true,
  breaklines=true,                              % parte las líneas largas
  breakatwhitespace=false,
  postbreak=\mbox{\textcolor{codigoGris}{$\hookrightarrow$}\space},
  numbers=left,
  numberstyle=\tiny\color{codigoGris},
  numbersep=8pt,
  xleftmargin=17pt,
  keywordstyle=\color{codigoClave}\bfseries,
  stringstyle=\color{codigoCadena},
  commentstyle=\color{codigoComentario}\itshape,
  emphstyle=\color{codigoTipo},
  captionpos=b,
  literate=%
    {á}{{\'a}}1 {é}{{\'e}}1 {í}{{\'i}}1 {ó}{{\'o}}1 {ú}{{\'u}}1
    {Á}{{\'A}}1 {É}{{\'E}}1 {Í}{{\'I}}1 {Ó}{{\'O}}1 {Ú}{{\'U}}1
    {ü}{{\"u}}1 {Ü}{{\"U}}1 {ñ}{{\~n}}1 {Ñ}{{\~N}}1
    {¿}{{?`}}1 {¡}{{!`}}1
    {«}{{\guillemotleft}}1 {»}{{\guillemotright}}1
    {—}{{---}}1 {–}{{--}}1 {…}{{\ldots}}1
    {º}{{\textordmasculine}}1 {ª}{{\textordfeminine}}1 {°}{{\textdegree}}1
    {→}{{$\rightarrow$}}1 {≈}{{$\approx$}}1 {∈}{{$\in$}}1
}

%--------------------------------------------------------------
%  LENGUAJES
%  Python, bash y SQL vienen con listings; JSON y el pseudocódigo
%  en español se definen aquí porque listings no los trae.
%--------------------------------------------------------------
\lstdefinestyle{python}{language=Python}

\lstdefinestyle{bash}{language=bash}

\lstdefinestyle{sql}{language=SQL}

\lstdefinelanguage{json}{
  alsoletter={-},
  morestring=[b]",
  morekeywords={true,false,null},
  sensitive=true,
  morecomment=[l]{//},
}

\lstdefinestyle{json}{language=json}

\lstdefinelanguage{seudo}{
  sensitive=false,
  morekeywords={algoritmo,entrada,salida,para,cada,hacer,fin,mientras,si,entonces,
    sino,devolver,retornar,repetir,hasta,funcion,procedimiento,inicio,
    verdadero,falso,no,y,o,listas,consulta,consulta,documento,documentos,
    ordenar,mezclar,fusionar,rango,puntuacion,puntuaciones},
  morecomment=[l]{\#},
}

\lstdefinestyle{seudocodigo}{language=seudo}

% Salida de consola: sin números de línea ni resaltado, para pegarla tal cual.
\lstdefinestyle{salida}{
  numbers=none,
  xleftmargin=4pt,
  keywordstyle={},
  stringstyle={},
  commentstyle=\color{codigoGris},
}

%--------------------------------------------------------------
%  CAJAS (tcolorbox)
%--------------------------------------------------------------
\tcbset{
  codigoCaja/.style={
    listing engine=listings,
    listing only,
    colback=codigoFondo,
    colframe=codigoLinea,
    coltitle=black,
    colbacktitle=codigoTitulo,
    fonttitle=\bfseries\small,
    boxrule=0.4pt,
    arc=1.5pt,
    outer arc=1.5pt,
    left=4pt, right=4pt, top=4pt, bottom=4pt,
    before skip=8pt, after skip=8pt,
  },
  % Alias cortos: \begin{codigo}[py]{…} equivale a [listing style=python].
  % La forma canónica de tcolorbox, [listing style=python] o
  % [listing options={…}], también funciona y sirve para cualquier opción
  % de listings (numbers=none, firstnumber=5, …).
  py/.style={listing style=python},
  sh/.style={listing style=bash},
  sql/.style={listing style=sql},
  js/.style={listing style=json},
  seudo/.style={listing style=seudocodigo},
  sal/.style={listing style=salida},
  % Rango de líneas de un archivo: \archivocodigo[lineas={10}{25}]{ruta}{Título}
  lineas/.style 2 args={listing options app={firstline=#1, lastline=#2}},
}

%--------------------------------------------------------------
%  ENTORNOS NUMERADOS
%  Los tres comparten el contador «listadocodigo», que se reinicia en
%  cada capítulo: el título sale como «Listado 4.1 · …».
%--------------------------------------------------------------
\newtcblisting[auto counter, number within=chapter, list inside=lol, list type=lstlisting]%
  {codigo}[2][]{%
  codigoCaja,
  breakable,                                   % se puede partir entre páginas
  list text={#2},                              % el «Índice de listados» muestra número y título
  title={Listado~\thetcbcounter\ \textperiodcentered\ #2},
  #1
}

\newtcblisting[use counter from=codigo, list inside=lol, list type=lstlisting]%
  {codigoflotante}[2][]{%
  codigoCaja,
  float, floatplacement=tbp,
  list text={#2},
  title={Listado~\thetcbcounter\ \textperiodcentered\ #2},
  #1
}

\newtcbinputlisting[use counter from=codigo, list inside=lol, list type=lstlisting]%
  {\archivocodigo}[3][]{%
  codigoCaja,
  breakable,
  list text={#3},
  title={Listado~\thetcbcounter\ \textperiodcentered\ #3},
  listing file={#2},
  #1
}

%--------------------------------------------------------------
%  CÓDIGO EN LÍNEA
%  Se definen SIN argumento a propósito: \lstinline lee el texto en modo
%  verbatim del flujo de entrada, y con el código como argumento de macro
%  falla («lstinline ended by EOL»). El delimitador es el carácter que
%  sigue a la llamada:
%      \py{revision="1.2"}     \sh|ssh iimas|     \cod{top-100}
%  Ojo: el contenido no puede llevar el carácter delimitador; para casos
%  raros, usa \verb|...| .
%--------------------------------------------------------------
\newcommand{\py}{\lstinline[style=python]}
\newcommand{\sh}{\lstinline[style=bash]}
\newcommand{\cod}{\lstinline}

%--------------------------------------------------------------
%  NOMBRES EN ESPAÑOL
%  Se fijan en \AtBeginDocument para que manden sobre los de babel.
%--------------------------------------------------------------
\AtBeginDocument{%
  \renewcommand{\lstlistingname}{Listado}%
  \renewcommand{\lstlistlistingname}{Índice de listados}%
}

%--------------------------------------------------------------
%  REFERENCIAS
%  Cualquiera de los tres entornos acepta la clave «label»:
%      \begin{codigoflotante}[py, label={lst:carga}]{Carga del split}
%  Así, \ref{lst:carga} imprime «5.1»; \autoref{lst:carga} (hyperref)
%  imprime «Listado 5.1»; y \cref{lst:carga} (cleveref) «listado 5.1».
%  cleveref es opcional: si no se carga, nada de esto estorba.
%  Dos cosas que costaron encontrar y conviene no volver a romper:
%   - El contador que crea tcolorbox se llama tcb@cnt@codigo. Es el nombre
%     con el que hay que registrar los nombres en español.
%   - Hay que registrarlos AL CARGAR cleveref, no en \AtBeginDocument:
%     cleveref construye los formatos de referencia al terminar de cargarse,
%     de modo que una definición posterior llega tarde y \cref imprime «??»
%     (aviso «reference format for label type … undefined»). Para eso está
%     el gancho \AddToHook{package/cleveref/after} de abajo.
%--------------------------------------------------------------
\AddToHook{package/cleveref/after}{%
  \Crefname{tcb@cnt@codigo}{Listado}{Listados}%
  \crefname{tcb@cnt@codigo}{listado}{listados}%
}

\makeatletter
\AtBeginDocument{%
  \expandafter\newcommand\csname tcb@cnt@codigoautorefname\endcsname{Listado}%
  %--- ÍNDICE DE LISTADOS EN EL ÍNDICE GENERAL
  % tocbibind solo conoce las listas de figuras y de tablas, así que el
  % «Índice de listados» no aparecía en el índice general. Se reusa su
  % propio mecanismo (\tocfile), que pone el encabezado y DESPUÉS el
  % \addcontentsline: así la entrada apunta a la página del encabezado
  % aunque la lista ocupe varias páginas.
  \@ifpackageloaded{tocbibind}{%
    \renewcommand{\lstlistoflistings}{\tocfile{\lstlistlistingname}{lol}}%
  }{}%
}
\makeatother

%  Documentación y ejemplos: notes/codigo.md  y  notes/codigo-ejemplos.tex
%
%  Resumen de la interfaz:
%    \begin{codigo}[style=python]{Título}   ... \end{codigo}          (se parte entre páginas)
%    \begin{codigoflotante}[style=python]{Título} ... \end{codigoflotante}  (flotante)
%    \archivocodigo[firstline=1,lastline=20]{ruta.py}{Título}         (desde archivo)
%    \py{código}   \sh{comando}   \cod{texto}                         (en línea)
%==============================================================

\usepackage{xcolor}          % Modelos de color RGB (ya lo carga tcolorbox, aquí para dejarlo explícito)
\usepackage{listings}        % Resaltado de sintaxis
\usepackage{upquote}         % Comillas rectas ' en el código (no tipográficas)
\usepackage{listingsutf8}    % Archivos completos en UTF-8 con \lstinputlisting
\usepackage[listings,breakable,skins,hooks]{tcolorbox}   % Marco, título, partido entre páginas

%--------------------------------------------------------------
%  PALETA
%  Sobria a propósito: la tesis se imprime en blanco y negro, así que
%  todo debe seguir siendo distinguible por contraste y por estilo.
%--------------------------------------------------------------
\definecolor{codigoFondo}{RGB}{252,252,250}      % fondo del listado
\definecolor{codigoTitulo}{RGB}{238,238,232}     % barra del título
\definecolor{codigoLinea}{RGB}{130,130,130}      % marco
\definecolor{codigoGris}{RGB}{95,95,95}          % números de línea y flecha de continuación
\definecolor{codigoClave}{RGB}{0,55,125}         % palabras clave
\definecolor{codigoCadena}{RGB}{140,35,35}       % cadenas
\definecolor{codigoComentario}{RGB}{85,105,85}   % comentarios
\definecolor{codigoTipo}{RGB}{20,85,105}         % tipos y clases

%--------------------------------------------------------------
%  CONFIGURACIÓN COMÚN
%  El mapa «literate» es lo que permite escribir acentos, eñes y
%  comillas latinas DENTRO de un listado con pdflatex: listings no
%  entiende UTF-8 multibyte por sí solo.
%--------------------------------------------------------------
\lstset{
  basicstyle=\ttfamily\small,
  columns=fullflexible,
  keepspaces=true,
  showstringspaces=false,
  tabsize=4,
  upquote=true,
  breaklines=true,                              % parte las líneas largas
  breakatwhitespace=false,
  postbreak=\mbox{\textcolor{codigoGris}{$\hookrightarrow$}\space},
  numbers=left,
  numberstyle=\tiny\color{codigoGris},
  numbersep=8pt,
  xleftmargin=17pt,
  keywordstyle=\color{codigoClave}\bfseries,
  stringstyle=\color{codigoCadena},
  commentstyle=\color{codigoComentario}\itshape,
  emphstyle=\color{codigoTipo},
  captionpos=b,
  literate=%
    {á}{{\'a}}1 {é}{{\'e}}1 {í}{{\'i}}1 {ó}{{\'o}}1 {ú}{{\'u}}1
    {Á}{{\'A}}1 {É}{{\'E}}1 {Í}{{\'I}}1 {Ó}{{\'O}}1 {Ú}{{\'U}}1
    {ü}{{\"u}}1 {Ü}{{\"U}}1 {ñ}{{\~n}}1 {Ñ}{{\~N}}1
    {¿}{{?`}}1 {¡}{{!`}}1
    {«}{{\guillemotleft}}1 {»}{{\guillemotright}}1
    {—}{{---}}1 {–}{{--}}1 {…}{{\ldots}}1
    {º}{{\textordmasculine}}1 {ª}{{\textordfeminine}}1 {°}{{\textdegree}}1
    {→}{{$\rightarrow$}}1 {≈}{{$\approx$}}1 {∈}{{$\in$}}1
}

%--------------------------------------------------------------
%  LENGUAJES
%  Python, bash y SQL vienen con listings; JSON y el pseudocódigo
%  en español se definen aquí porque listings no los trae.
%--------------------------------------------------------------
\lstdefinestyle{python}{language=Python}

\lstdefinestyle{bash}{language=bash}

\lstdefinestyle{sql}{language=SQL}

\lstdefinelanguage{json}{
  alsoletter={-},
  morestring=[b]",
  morekeywords={true,false,null},
  sensitive=true,
  morecomment=[l]{//},
}

\lstdefinestyle{json}{language=json}

\lstdefinelanguage{seudo}{
  sensitive=false,
  morekeywords={algoritmo,entrada,salida,para,cada,hacer,fin,mientras,si,entonces,
    sino,devolver,retornar,repetir,hasta,funcion,procedimiento,inicio,
    verdadero,falso,no,y,o,listas,consulta,consulta,documento,documentos,
    ordenar,mezclar,fusionar,rango,puntuacion,puntuaciones},
  morecomment=[l]{\#},
}

\lstdefinestyle{seudocodigo}{language=seudo}

% Salida de consola: sin números de línea ni resaltado, para pegarla tal cual.
\lstdefinestyle{salida}{
  numbers=none,
  xleftmargin=4pt,
  keywordstyle={},
  stringstyle={},
  commentstyle=\color{codigoGris},
}

%--------------------------------------------------------------
%  CAJAS (tcolorbox)
%--------------------------------------------------------------
\tcbset{
  codigoCaja/.style={
    listing engine=listings,
    listing only,
    colback=codigoFondo,
    colframe=codigoLinea,
    coltitle=black,
    colbacktitle=codigoTitulo,
    fonttitle=\bfseries\small,
    boxrule=0.4pt,
    arc=1.5pt,
    outer arc=1.5pt,
    left=4pt, right=4pt, top=4pt, bottom=4pt,
    before skip=8pt, after skip=8pt,
  },
  % Alias cortos: \begin{codigo}[py]{…} equivale a [listing style=python].
  % La forma canónica de tcolorbox, [listing style=python] o
  % [listing options={…}], también funciona y sirve para cualquier opción
  % de listings (numbers=none, firstnumber=5, …).
  py/.style={listing style=python},
  sh/.style={listing style=bash},
  sql/.style={listing style=sql},
  js/.style={listing style=json},
  seudo/.style={listing style=seudocodigo},
  sal/.style={listing style=salida},
  % Rango de líneas de un archivo: \archivocodigo[lineas={10}{25}]{ruta}{Título}
  lineas/.style 2 args={listing options app={firstline=#1, lastline=#2}},
}

%--------------------------------------------------------------
%  ENTORNOS NUMERADOS
%  Los tres comparten el contador «listadocodigo», que se reinicia en
%  cada capítulo: el título sale como «Listado 4.1 · …».
%--------------------------------------------------------------
\newtcblisting[auto counter, number within=chapter, list inside=lol, list type=lstlisting]%
  {codigo}[2][]{%
  codigoCaja,
  breakable,                                   % se puede partir entre páginas
  list text={#2},                              % el «Índice de listados» muestra número y título
  title={Listado~\thetcbcounter\ \textperiodcentered\ #2},
  #1
}

\newtcblisting[use counter from=codigo, list inside=lol, list type=lstlisting]%
  {codigoflotante}[2][]{%
  codigoCaja,
  float, floatplacement=tbp,
  list text={#2},
  title={Listado~\thetcbcounter\ \textperiodcentered\ #2},
  #1
}

\newtcbinputlisting[use counter from=codigo, list inside=lol, list type=lstlisting]%
  {\archivocodigo}[3][]{%
  codigoCaja,
  breakable,
  list text={#3},
  title={Listado~\thetcbcounter\ \textperiodcentered\ #3},
  listing file={#2},
  #1
}

%--------------------------------------------------------------
%  CÓDIGO EN LÍNEA
%  Se definen SIN argumento a propósito: \lstinline lee el texto en modo
%  verbatim del flujo de entrada, y con el código como argumento de macro
%  falla («lstinline ended by EOL»). El delimitador es el carácter que
%  sigue a la llamada:
%      \py{revision="1.2"}     \sh|ssh iimas|     \cod{top-100}
%  Ojo: el contenido no puede llevar el carácter delimitador; para casos
%  raros, usa \verb|...| .
%--------------------------------------------------------------
\newcommand{\py}{\lstinline[style=python]}
\newcommand{\sh}{\lstinline[style=bash]}
\newcommand{\cod}{\lstinline}

%--------------------------------------------------------------
%  NOMBRES EN ESPAÑOL
%  Se fijan en \AtBeginDocument para que manden sobre los de babel.
%--------------------------------------------------------------
\AtBeginDocument{%
  \renewcommand{\lstlistingname}{Listado}%
  \renewcommand{\lstlistlistingname}{Índice de listados}%
}

%--------------------------------------------------------------
%  REFERENCIAS
%  Cualquiera de los tres entornos acepta la clave «label»:
%      \begin{codigoflotante}[py, label={lst:carga}]{Carga del split}
%  Así, \ref{lst:carga} imprime «5.1»; \autoref{lst:carga} (hyperref)
%  imprime «Listado 5.1»; y \cref{lst:carga} (cleveref) «listado 5.1».
%  cleveref es opcional: si no se carga, nada de esto estorba.
%  Dos cosas que costaron encontrar y conviene no volver a romper:
%   - El contador que crea tcolorbox se llama tcb@cnt@codigo. Es el nombre
%     con el que hay que registrar los nombres en español.
%   - Hay que registrarlos AL CARGAR cleveref, no en \AtBeginDocument:
%     cleveref construye los formatos de referencia al terminar de cargarse,
%     de modo que una definición posterior llega tarde y \cref imprime «??»
%     (aviso «reference format for label type … undefined»). Para eso está
%     el gancho \AddToHook{package/cleveref/after} de abajo.
%--------------------------------------------------------------
\AddToHook{package/cleveref/after}{%
  \Crefname{tcb@cnt@codigo}{Listado}{Listados}%
  \crefname{tcb@cnt@codigo}{listado}{listados}%
}

\makeatletter
\AtBeginDocument{%
  \expandafter\newcommand\csname tcb@cnt@codigoautorefname\endcsname{Listado}%
  %--- ÍNDICE DE LISTADOS EN EL ÍNDICE GENERAL
  % tocbibind solo conoce las listas de figuras y de tablas, así que el
  % «Índice de listados» no aparecía en el índice general. Se reusa su
  % propio mecanismo (\tocfile), que pone el encabezado y DESPUÉS el
  % \addcontentsline: así la entrada apunta a la página del encabezado
  % aunque la lista ocupe varias páginas.
  \@ifpackageloaded{tocbibind}{%
    \renewcommand{\lstlistoflistings}{\tocfile{\lstlistlistingname}{lol}}%
  }{}%
}
\makeatother
